\ifdefined\OTCOMBINED%
\else
\documentclass{otscience}
% ============================================================
% Shared setup for OT Science modular workbook parts
% Place these files in the same folder as your fixed otscience.cls
% ============================================================

\makeatletter
\makeatother
\usepackage{multicol}
\usepackage{longtable}
\usepackage{makecell}
\usepackage{pdflscape}
\usepackage{bm}
\providecommand{\blank}[1][3cm]{\rule{#1}{0.4pt}}
\providecommand{\bigblank}{\rule{7cm}{0.4pt}}
\providecommand{\componentblank}{\blank[2.5cm]}
\providecommand{\R}{\mathbb{R}}
\providecommand{\C}{\mathbb{C}}
\providecommand{\ii}{\mathrm{i}}
\providecommand{\ee}{\mathrm{e}}
\providecommand{\dd}{\mathrm{d}}
\providecommand{\grad}{\nabla}
\providecommand{\divergence}{\nabla\!\cdot}
\providecommand{\curl}{\nabla\!\times}
\providecommand{\lap}{\nabla^2}
\providecommand{\ihat}{\hat{\mathbf{i}}}
\providecommand{\jhat}{\hat{\mathbf{j}}}
\providecommand{\khat}{\hat{\mathbf{k}}}
\providecommand{\er}{\hat{\mathbf{e}}_r}
\providecommand{\etheta}{\hat{\mathbf{e}}_\theta}
\providecommand{\ephi}{\hat{\mathbf{e}}_\phi}
\providecommand{\erho}{\hat{\boldsymbol{\rho}}}
\providecommand{\vA}{\mathbf{A}}
\providecommand{\vB}{\mathbf{B}}
\providecommand{\va}{\mathbf{a}}
\providecommand{\vb}{\mathbf{b}}
\providecommand{\vc}{\mathbf{c}}
\providecommand{\vx}{\mathbf{x}}
\providecommand{\vr}{\mathbf{r}}
\providecommand{\matA}{\mathbf{A}}
\providecommand{\matB}{\mathbf{B}}
\providecommand{\tr}{\operatorname{tr}}
\providecommand{\cof}{\operatorname{cof}}
\providecommand{\dev}{\operatorname{dev}}
\providecommand{\diag}{\operatorname{diag}}
\providecommand{\questionline}[1][5cm]{\rule{#1}{0.4pt}}
\setlength{\parindent}{0pt}
\setlength{\parskip}{4pt}
\renewcommand{\arraystretch}{1.25}

% Fallbacks in case the current otscience.cls has not defined nosplit boxes.
\makeatletter
\@ifundefined{formulaboxnosplit}{\newenvironment{formulaboxnosplit}[1][Formula]{\begin{formulabox}[#1]}{\end{formulabox}}}{}
\@ifundefined{definitionboxnosplit}{\newenvironment{definitionboxnosplit}[1][Definition]{\begin{definitionbox}[#1]}{\end{definitionbox}}}{}
\@ifundefined{summaryboxnosplit}{\newenvironment{summaryboxnosplit}[1][Summary]{\begin{summarybox}[#1]}{\end{summarybox}}}{}
\@ifundefined{warningboxnosplit}{\newenvironment{warningboxnosplit}[1][Warning]{\begin{warningbox}[#1]}{\end{warningbox}}}{}
\@ifundefined{exampleboxnosplit}{\newenvironment{exampleboxnosplit}[1][Example]{\begin{examplebox}[#1]}{\end{examplebox}}}{}
\makeatother

\newenvironment{practicebox}[1][Quick Drills and Practice]{\begin{examplebox}[#1]}{\end{examplebox}}
\newenvironment{practiceboxnosplit}[1][Quick Drills and Practice]{\begin{exampleboxnosplit}[#1]}{\end{exampleboxnosplit}}

\newcommand{\standalonetitle}[3]{%
    \title{\textbf{#1}\\\large #2}
    \author{OT Science Notes}
    \date{#3}
}

\standalonetitle{Index Notation, Kronecker Delta and Levi-Civita Symbol}{Part 3 of the OT Science Formula Completion Workbook}{\DTMtoday}
\begin{document}
\maketitle
\tableofcontents
\newpage
\fi

% 03_index_delta_levicivita.tex
\section{Index Notation, Kronecker Delta and Levi-Civita Symbol}
\sectionmark{Index Notation}

Index notation compresses vector and tensor formulas into short expressions. It is especially useful when proving vector identities. The price is that every index must be handled carefully.

\subsection{Einstein Summation Convention}

When an index is repeated once as a subscript or repeated in a product in Euclidean notation, it is summed over its full range. In three-dimensional Cartesian coordinates, the range is usually \(1,2,3\).

\begin{formulaboxnosplit}[Einstein Summation]
    \[
        a_i b_i=a_1b_1+a_2b_2+a_3b_3.
    \]

    \[
        A_{ij}b_j=A_{i1}b_1+A_{i2}b_2+A_{i3}b_3.
    \]

    \[
        A_{ij}B_{ij}=\sum_{i=1}^{3}\sum_{j=1}^{3}A_{ij}B_{ij}.
    \]

    \[
        {(A B)}_{ij}=A_{ik}B_{kj}.
    \]
\end{formulaboxnosplit}

A free index appears on both sides of an equation and is not summed. A dummy index is repeated and summed. Dummy indices can be renamed as long as no conflict is created.

\subsection{Kronecker Delta}

The Kronecker delta is the identity tensor in component form. It replaces one index by another.

\begin{formulabox}[Kronecker Delta Identities]
    \[
        \delta_{ij}=\begin{cases}
            1, & i=j,     \\
            0, & i\neq j.
        \end{cases}
    \]

    \[
        \delta_{ij}a_j=a_i,
        \qquad
        \delta_{ij}a_i=a_j.
    \]

    \[
        \delta_{ij}\delta_{jk}=\delta_{ik},
        \qquad
        \delta_{ii}=3.
    \]

    \[
        \delta_{ij}A_{ij}=A_{ii}=\tr(A).
    \]

    \[
        \partial_i x_j=\frac{\partial x_j}{\partial x_i}=\delta_{ij}.
    \]

    \[
        \partial_i x_i=3,
        \qquad
        \partial_i r=\frac{x_i}{r},
        \qquad
        \partial_i(r^2)=2x_i.
    \]
\end{formulabox}

\subsection{Levi-Civita Symbol}

The Levi-Civita symbol tracks orientation and cross products. It is zero whenever any two indices are repeated.

\begin{formulabox}[Levi-Civita Identities]
    \[
        \varepsilon_{ijk}=\begin{cases}
            +1, & (i,j,k)\text{ is an even permutation of }(1,2,3), \\
            -1, & (i,j,k)\text{ is an odd permutation of }(1,2,3),  \\
            0,  & \text{two indices are repeated}.
        \end{cases}
    \]

    \[
        {(\mathbf{a}\times\mathbf{b})}_i=\varepsilon_{ijk}a_{j}b_k.
    \]

    \[
        \varepsilon_{ijk}\varepsilon_{ilm}=\delta_{jl}\delta_{km}-\delta_{jm}\delta_{kl}.
    \]

    \[
        \varepsilon_{ijk}\varepsilon_{ijm}=2\delta_{km},
        \qquad
        \varepsilon_{ijk}\varepsilon_{ijk}=6.
    \]

    \[
        \varepsilon_{ijk}\varepsilon_{lmn}=\begin{vmatrix}
            \delta_{il} & \delta_{im} & \delta_{in} \\
            \delta_{jl} & \delta_{jm} & \delta_{jn} \\
            \delta_{kl} & \delta_{km} & \delta_{kn}
        \end{vmatrix}.
    \]
\end{formulabox}

\subsection{Vector Calculus in Index Notation}

\begin{formulabox}[Vector Operations in Index Form]
    \[
        {(\nabla f)}_i=\partial_i f.
    \]

    \[
        \nabla\cdot\mathbf{A}=\partial_{i}A_i.
    \]

    \[
        {(\nabla\times\mathbf{A})}_i=\varepsilon_{ijk}\partial_{j}A_k.
    \]

    \[
        \nabla^2 f=\partial_i\partial_i f.
    \]

    \[
        {(\nabla\mathbf{A})}_{ij}=\partial_{j}A_i.
    \]

    \[
        {(\nabla\cdot\mathbf{T})}_i=\partial_{j}T_{ij}.
    \]
\end{formulabox}

\subsection{Proof Pattern: Curl of Curl}

The index method is powerful because the proof becomes a controlled substitution. Starting with
\[
    {(\nabla\times(\nabla\times\mathbf{A}))}_i=\varepsilon_{ijk}\partial_j{(\nabla\times\mathbf{A})}_k,
\]
use
\[
    {(\nabla\times\mathbf{A})}_k=\varepsilon_{klm}\partial_{l}A_m.
\]
Then
\[
    {(\nabla\times(\nabla\times\mathbf{A}))}_i=\varepsilon_{ijk}\varepsilon_{klm}\partial_j\partial_{l}A_m.
\]
Using the contraction identity gives the familiar result
\[
    \nabla\times(\nabla\times\mathbf{A})=\nabla(\nabla\cdot\mathbf{A})-\nabla^2\mathbf{A}.
\]

\begin{warningboxnosplit}[Index-Notation Mistakes]
    \begin{enumerate}
        \item Do not repeat the same index more than twice in one term unless you have explicitly defined what it means.
        \item Do not change a free index on one side of an equation without changing it on the other side.
        \item Dummy indices may be renamed, but free indices may not.
        \item The identity \(\varepsilon_{ijk}\varepsilon_{ilm}\) depends on which index is contracted. Check the repeated index carefully.
    \end{enumerate}
\end{warningboxnosplit}

\begin{practicebox}[Quick Drills and Practice: Index Notation]
    \textbf{Level A:\,Expansion and simplification}
    \begin{multicols}{2}
        \begin{enumerate}
            \item Expand \(a_{i}b_i\) in three dimensions.
            \item Expand \(A_{ij}b_j\) for \(i=1\).
            \item Simplify \(\delta_{ij}a_j\).
            \item Simplify \(\delta_{ij}A_{jk}\).
            \item Simplify \(\delta_{ij}\delta_{jk}\delta_{kl}\).
            \item Simplify \(\delta_{ij}A_{ij}\).
            \item Compute \(\delta_{ii}\) in three dimensions.
            \item Compute \(\partial_i x_j\).
            \item Compute \(\partial_i x_i\).
            \item Compute \(\partial_i(x_{i}x_j)\).
        \end{enumerate}
    \end{multicols}

    \textbf{Level B:\,Levi-Civita}
    \begin{enumerate}
        \item Evaluate \(\varepsilon_{123}\), \(\varepsilon_{132}\), \(\varepsilon_{221}\), \(\varepsilon_{312}\).
        \item Expand \({(\mathbf{a}\times\mathbf{b})}_1=\varepsilon_{1jk}a_{j}b_k\).
        \item Show that \(\mathbf{a}\times\mathbf{b}=-\mathbf{b}\times\mathbf{a}\) using \(\varepsilon_{ijk}\).
        \item Simplify \(\varepsilon_{ijk}\varepsilon_{ijm}\).
        \item Simplify \(\varepsilon_{ijk}\varepsilon_{ijk}\).
        \item Use \(\varepsilon_{ijk}\varepsilon_{ilm}\) to simplify \((\mathbf{a}\times\mathbf{b})\cdot(\mathbf{c}\times\mathbf{d})\).
    \end{enumerate}

    \textbf{Level C:\,Vector identities}
    \begin{enumerate}
        \item Prove \(\nabla\cdot(\nabla\times\mathbf{A})=0\) using index notation.
        \item Prove \(\nabla\times(\nabla f)=\mathbf{0}\) using index notation.
        \item Prove \(\nabla\times(\nabla\times\mathbf{A})=\nabla(\nabla\cdot\mathbf{A})-\nabla^2\mathbf{A}\).
        \item Prove \(\mathbf{a}\times(\mathbf{b}\times\mathbf{c})=\mathbf{b}(\mathbf{a}\cdot\mathbf{c})-\mathbf{c}(\mathbf{a}\cdot\mathbf{b})\).
        \item Prove \((\mathbf{a}\times\mathbf{b})\cdot(\mathbf{c}\times\mathbf{d})=(\mathbf{a}\cdot\mathbf{c})(\mathbf{b}\cdot\mathbf{d})-(\mathbf{a}\cdot\mathbf{d})(\mathbf{b}\cdot\mathbf{c})\).
    \end{enumerate}

    \textbf{Level D:\,Derivative identities}
    \begin{enumerate}
        \item Show that \(\partial_i r=x_i/r\), where \(r={(x_{j}x_j)}^{1/2}\).
        \item Show that \(\partial_i(1/r)=-x_i/r^3\).
        \item Compute \(\partial_i(x_i/r^3)\) for \(r\neq0\).
        \item Compute \(\partial_j(x_{i}x_j)\) in three dimensions.
        \item If \(T_{ij}=f\delta_{ij}\), compute \(\partial_{j}T_{ij}\).
    \end{enumerate}
\end{practicebox}

\ifdefined\OTCOMBINED%
\else
\end{document}
\fi
